\documentclass[10pt,a4paper]{amsart}
\usepackage[margin=19mm]{geometry}
\usepackage{amsmath,amssymb,mathtools}
\usepackage[scr=rsfs,cal=euler]{mathalfa}
\usepackage{xfrac}
\usepackage[unicode,hidelinks,bookmarks=false]{hyperref}
\newcommand{\ie}{i.e.\ }
\newcommand{\cf}{cf.\ }
\newcommand{\resp}{respectively\ }
\newcommand{\Subsection}{\subsection}
\providecommand{\nobreakdash}{\nobreak\hbox{-}}
\setlength{\emergencystretch}{2em}
\multlinegap0pt
\usepackage{bm}
\usepackage{bbm}
\usepackage{dsfont}
\usepackage{xspace}
\usepackage{cancel}
\usepackage{mathtools}
\usepackage{amsthm}
\makeatletter
\@ifundefined{theorem}{\newtheorem{theorem}{Theorem}}{}
\@ifundefined{lemma}{\newtheorem{lemma}[theorem]{Lemma}}{}
\makeatother
\title[2-Elongated Plane Partitions and Powers of 7: The Localizati]{2-Elongated Plane Partitions and Powers of 7: The Localization Method Applied to a Genus 1 Congruence Family}
\author{Koustav Banerjee, Nicolas Allen Smoot}
\date{Source: arXiv:2306.15594v4 (CC BY 4.0)}
\begin{document}
\maketitle

\noindent\emph{Source and licence.} 2-Elongated Plane Partitions and Powers of 7: The Localization Method Applied to a Genus 1 Congruence Family. Version arXiv:2306.15594v4 (CC BY 4.0), distributed under the Creative Commons Attribution 4.0 International licence, as recorded in the arXiv licence metadata for this exact version and shown on the abstract page https://arxiv.org/abs/2306.15594v4. Full rights notice: licenses/arxiv-2306.15594-CC-BY-4.0.txt.\par\smallskip

\noindent\textit{Editorial scope.} Counted unit: the Lemma whose source label is \texttt{lowerboundpiabc} in the article (source0.tex lines 458--466, source label \texttt{lowerboundpiabc}), delivered together with the article's own complete original proof of it (source0.tex lines 468--522, one \texttt{proof} environment of 55 lines). No mathematics is written, repaired or re-derived: statement and proof are the article's own text line for line. Required context retained uncounted: ufgreduceA (lines 358--360); modeqnz (lines 410--415); modeqnx (lines 421--423). Every definition, display and label of those lines is retained unchanged. Omitted from this package and not redistributed: 1--357, 361--409, 416--420, 424--457, 467--467, 523--1619. Nothing inside a retained excerpt is omitted or altered: every retained body is delivered exactly as the article's lines are written, so no omission receipt of any kind accompanies this package. Citations: the retained excerpts carry no citation command at all, so no bibliography entry of the article is transcribed and no reference is redistributed. Reference adaptation: none was needed and none was made; every reference inside the delivered unit resolves inside it, so no unresolved reference or citation is produced. Printed slips kept literal: no printed word, symbol, hypothesis or number of the counted statement or of its proof was altered. Layout and class: the article's own class is \texttt{amsart} with packages including url, algorithmic, geometry; the wrapper is the lane's pinned \texttt{amsart} editorial wrapper at 10pt on A4, which restates the theorem-like environments the retained text needs and adds the article's own macro definitions that the retained text needs (none), with extra wrapper packages (bm, bbm, dsfont, xspace, cancel, mathtools). The delivered numbering is the unit's own. Assets: no retained span contains a figure environment, a graphics-inclusion command, a PDF-page inclusion, a table or a TikZ picture; no image file is copied into this package, the package declares no asset dependency and no image file is redistributed with it. Licence: version arXiv:2306.15594v4 of this article is distributed under the Creative Commons Attribution 4.0 International licence, as recorded in the arXiv licence metadata for this exact version and shown on the abstract page https://arxiv.org/abs/2306.15594v4. A full rights notice is delivered at licenses/arxiv-2306.15594-CC-BY-4.0.txt. Characters: every retained excerpt is pure ASCII, so the delivered engine, XeLaTeX with its default Latin Modern text font, reports no missing character.
\begin{align}
U_{\ell}\left(f(q^{\ell})g(q)\right) = f(q) U_{\ell}\left(g(q)\right).\label{ufgreduceA}
\end{align}

\begin{theorem}\label{modeqnz}
For $0\le k\le 14$, let $b_k\in\mathbb{Z}[X]$ be defined as in the Appendix.  We have
\begin{align}
z^{14}+\sum_{k=0}^{13}b_k(z(7\tau))z^k=0.\label{actmodeqnz}
\end{align}
\end{theorem}  Notice that by (\ref{ufgreduceA}), $U_{7}\left( b_k(z(7\tau)) \right)=b_k(z(\tau))$.

\begin{align}
x^{14}+\sum_{j=0}^{13}a_j(x(7\tau))x^j=0.\label{modeqnx}
\end{align}

\begin{lemma}\label{lowerboundpiabc}
Suppose that we define 
\begin{align}
\hat{\pi}^{(\alpha)}_{\beta\gamma}(m,r) := \left\lfloor \frac{7r-m+\epsilon^{(\alpha)}_{\beta\gamma}}{9} \right\rfloor\label{pihatvalues}
\end{align} for some $\epsilon^{(\alpha)}_{\beta\gamma}\in\mathbb{Z}$.  If the relation
\begin{align}
U^{(\alpha)}\left( \frac{y^{\beta}x^{m}}{(1+7x)^{n}} \right) =& \frac{1}{(1+7x)^{7n+\kappa}}\sum_{\substack{0\le\gamma\le 1\\ r\ge 1-\gamma}} h^{(\alpha)}_{\beta\gamma}(m,n,r) 7^{\hat{\pi}^{(\alpha)}_{\beta\gamma}(m,r)}y^{\gamma}x^r\label{lowerbounduaymoz}
\end{align} holds for $1\le m\le 14$, $1\le n\le 14$, then it must also hold for all $m,n\ge 15$.
\end{lemma}

\begin{proof}

We begin by taking (\ref{modeqnz}):
\begin{align*}
z^{14}+\sum_{k=0}^{13}b_k(z(7\tau))z^k= \sum_{k=1}^{14}b_k(z(7\tau))z^k + b_0(z(7\tau)).
\end{align*}  Rearranging and dividing through by $b_0(z(7\tau))$, we have
\begin{align*}
1 =& \frac{-1}{b_0(z(7\tau))}\sum_{k=1}^{14}b_k(z(7\tau))\cdot z^{k}.
\end{align*}  We now divide both sides by $z^n$:
\begin{align*}
\frac{1}{z^{n}} =& \frac{-1}{z(7\tau)^{14}}\sum_{k=1}^{14}b_k(z(7\tau))\cdot \frac{1}{z^{n-k}}.
\end{align*}  We now multiply by a given power of $x$ and $y$:
\begin{align*}
\frac{y^{\beta}x^m}{z^{n}} =& \frac{-1}{z(7\tau)^{14}}\sum_{k=1}^{14}b_k(z(7\tau))\frac{y^{\beta}x^m}{z^{n-k}}.
\end{align*}  Applying $U^{(\alpha)}$ to both sides and and substituting $z=1+7x$, we have
\begin{align}
U^{(\alpha)}\left(\frac{y^{\beta}x^m}{(1+7x)^{n}}\right) =& \frac{-1}{(1+7x)^{14}}\sum_{k=1}^{14}b_k(z(\tau)) U^{(\alpha)}\left(\frac{y^{\beta}x^m}{(1+7x)^{n-k}}\right).\label{ualphaxmznmk}
\end{align}  Next, we recall (\ref{modeqnx}):
\begin{align*}
x^{14}+\sum_{j=0}^{13}a_j(x(7\tau))x^j=0.
\end{align*} 
If we multiply both sides by $\frac{y^{\beta}x^{m-14}}{z^{n-k}}$ for some $m\ge 14$, we have
\begin{align*}
\frac{y^{\beta}x^{m}}{z^{n-k}}+\sum_{j=0}^{13}a_j(x(7\tau))\frac{y^{\beta}x^{m+j-14}}{z^{n-k}}=0.
\end{align*}  Applying $U^{(\alpha)}$, rearranging, and again substituting $z=1+7x$, we have
\begin{align*}
U^{(\alpha)}\left(\frac{y^{\beta}x^{m}}{(1+7x)^{n-k}}\right) = -\sum_{j=0}^{13}a_j(x(\tau)) U^{(\alpha)}\left(\frac{y^{\beta}x^{m+j-14}}{(1+7x)^{n-k}}\right).
\end{align*}  Substituting this into (\ref{ualphaxmznmk}), we now also take $n\ge 14$.  We show that if (\ref{lowerbounduaymoz}) applies to $U^{(\alpha)}\left( \frac{y^{\beta}x^{m+j-14}}{(1+7x)^{n-k}} \right)$ for $0\le j\le 13$, $1\le k\le 14$, then the relation must apply to $U^{(\alpha)}\left( \frac{y^{\beta}x^{m}}{(1+7x)^{n}} \right)$.  Thus, if we verify the relation for the first 14 consecutive integral values of $m,n$, then the relation will apply for all higher $m,n$.  To this end, we have
\begin{align*}
U^{(\alpha)}&\left( \frac{y^{\beta}x^{m}}{(1+7x)^{n}} \right) = \frac{-1}{(1+7x)^{14}}\sum_{k=1}^{14}b_k(z(\tau))\cdot U^{(\alpha)}\left( \frac{y^{\beta}x^{m}}{(1+7x)^{n-k}} \right)\\
=& \frac{1}{(1+7x)^{14}}\sum_{j=0}^{13}\sum_{k=1}^{14}a_j(\tau)b_k(z(\tau))\cdot U^{(\alpha)}\left( \frac{y^{\beta}x^{m+j-14}}{(1+7x)^{n-k}} \right)\\
=& \frac{1}{(1+7x)^{14}}\sum_{j=0}^{13}\sum_{k=1}^{14}a_j(\tau)b_k(z(\tau))\cdot\frac{1}{(1+7x)^{7(n-k)+\kappa}}\sum_{\substack{0\le\gamma\le 1,\\ r\ge 1-\gamma}} h^{(\alpha)}_{\beta\gamma}(m+j-14,n-k,r) 7^{\pi^{(\alpha)}_{\beta\gamma}(m+j-14,r)}y^{\gamma}x^r\\
=& \frac{1}{(1+7x)^{7n+\kappa}}\sum_{j=0}^{13}\sum_{k=1}^{14}w(j,k) \sum_{\substack{0\le\gamma\le 1,\\ r\ge 1-\gamma}} h^{(\alpha)}_{\beta\gamma}(m+j-14,n-k,r) 7^{\pi^{(\alpha)}_{\beta\gamma}(m+j-14,r)}y^{\gamma}x^r,
\end{align*} where we define
\begin{align*}
w(j,k) :=& a_j(\tau)b_k(z(\tau))(1+7x)^{7(k-2)}\\
=& \sum_{l=1}^{L} v(j,k,l)\cdot 7^{\left\lfloor \frac{7l+j-6}{9} \right\rfloor}\cdot x^l.
\end{align*}  We verify that $w(j,k)$ has this form in our Mathematica supplement.  This gives us
\begin{align*}
U^{(\alpha)}&\left( \frac{y^{\beta}x^{m}}{(1+7x)^{n}} \right)\\
=& \frac{1}{(1+7x)^{7n+\kappa}}\sum_{\substack{0\le j\le 13,\\ 1\le k\le 14\\ 0\le l\le L\\0\le\gamma\le 1\\ r\ge 1-\gamma}}v(j,k,l) h^{(\alpha)}_{\beta\gamma}(m+j-14,n-k,r) 7^{\pi^{(\alpha)}_{\beta\gamma}(m+j-14,r)+\left\lfloor \frac{7l+j-6}{9} \right\rfloor}y^{\gamma}x^{r+l}\\
=& \frac{1}{(1+7x)^{7n+\kappa}}\sum_{\substack{0\le j\le 13,\\ 1\le k\le 14\\ 0\le l\le L\\0\le\gamma\le 1\\ r\ge 1+l-\gamma}}v(j,k,l) h^{(\alpha)}_{\beta\gamma}(m+j-14,n-k,r-l) 7^{\pi^{(\alpha)}_{\beta\gamma}(m+j-14,r-l)+\left\lfloor \frac{7l+j-6}{9} \right\rfloor}y^{\gamma}x^{r},
\end{align*} with the latter line coming from adjusting $r$.  Examining the powers of 7, we note that 
\begin{align*}
\pi^{(\alpha)}_{\beta\gamma}(m+j-14,r)+\left\lfloor \frac{7l+j-6}{9} \right\rfloor &= \left\lfloor \frac{7r-(m+j-14)+\epsilon^{(\alpha)}_{\beta\gamma}}{9} \right\rfloor +\left\lfloor \frac{7l+j-6}{9} \right\rfloor\\
&\ge \left\lfloor \frac{7(r+l)-m+\epsilon^{(\alpha)}_{\beta\gamma}}{9} \right\rfloor\\
&= \pi^{(\alpha)}_{\beta\gamma}(m,r+l).
\end{align*}  That is, for $r\ge l$,
\begin{align*}
\pi^{(\alpha)}_{\beta\gamma}(m+j-14,r-l)+\left\lfloor \frac{7l+j-6}{9} \right\rfloor &\ge \pi^{(\alpha)}_{\beta\gamma}(m,r).
\end{align*}  Therefore, the power of 7 necessary for (\ref{lowerbounduaymoz}) is achieved in the coefficient of $y^{\gamma}x^r$.  We now note that we can define
\begin{align*}
h^{(\alpha)}_{\beta\gamma}(m,n,r):=\sum_{\substack{0\le j\le 13,\\ 1\le k\le 14\\ 0\le l\le L}}v(j,k,l) h^{(\alpha)}_{\beta\gamma}(m+j-14,n-k,r-l) 7^{\pi^{(\alpha)}_{\beta\gamma}(m+j-14,r-l)+\left\lfloor \frac{7l+j-6}{9} \right\rfloor - \pi^{(\alpha)}_{\beta\gamma}(m,r)},
\end{align*} since the latter also has finite support in $r$.
\end{proof}

\end{document}
